\chapter{Centralne twierdzenie graniczne}
\label{ch:clt}

\providecommand{\atlasref}[2]{#2}

\paragraph{Odsyłacze międzytomowe.}
Zastosowania centralnego twierdzenia granicznego
do testowania hipotez przedstawiono w tomie IV,
w rozdziale o testowaniu%
\IfFileExists{probability_volume_iv.aux}{ (rozdział~\ref{IV-ch:testing})}{},
a metodę delty w rozdziale o asymptotycznej teorii estymacji%
\IfFileExists{probability_volume_iv.aux}{ (rozdział~\ref{IV-ch:asymptotic-estimation})}{}.


Prawa wielkich liczb opisują zbieżność średnich,
ale nie określają \emph{kształtu rozkładu} ich fluktuacji.
Dla parami niezależnych rzeczywistych zmiennych
losowych $X_1,X_2,\ldots$ na jednej przestrzeni
probabilistycznej, o wspólnej średniej $m$ i wspólnej
skończonej wariancji $\sigma^2$, oznaczmy
$\overline X_n=n^{-1}\sum_{k=1}^nX_k$.
Z rozdziału~\ref{ch:lln} wiemy, że
\[
\E|\overline X_n-m|^2=\Var(\overline X_n)=\frac{\sigma^2}{n}.
\]
Zatem $\|\overline X_n-m\|_{L^2(\Prob)}=\sigma/\sqrt n$,
gdzie $\sigma$ oznacza nieujemny pierwiastek z wariancji. Przy $\sigma>0$ naturalną skalą
odchylenia jest $n^{-1/2}$. Centralne twierdzenie
graniczne (CTG) wyjaśnia, pod jakimi dodatkowymi
założeniami otrzymamy rozkład normalny jako granicę
rozkładów odpowiednio przeskalowanych odchyleń.

\section{Centralne twierdzenie graniczne dla prób niezależnych jednakowo rozłożonych}

\begin{twierdzenie}[Centralne twierdzenie graniczne Lindeberga--Lévy'ego]
Niech $X_1,X_2,\ldots$ będą wzajemnie niezależnymi
rzeczywistymi zmiennymi losowymi na jednej przestrzeni
probabilistycznej, o tym samym rozkładzie, oraz niech
$\E X_1^2<\infty$ i $\Var(X_1)>0$.
Warunek $\E X_1^2<\infty$ zapewnia całkowalność $X_1$
oraz skończoność wariancji (rozdział~\ref{ch:expectation}).
Oznaczmy $m=\E X_1$, $\sigma=\sqrt{\Var(X_1)}>0$,
$S_n=\sum_{k=1}^nX_k$ oraz $\overline X_n=S_n/n$
dla $n\geq1$. Wówczas zachodzi
\[
\frac{S_n-nm}{\sigma\sqrt n}
=\frac{\sqrt n(\overline X_n-m)}{\sigma}
\xrightarrow{d}\Normal(0,1).
\]
Zapis z rozkładem po prawej stronie oznacza słabą
zbieżność miar obrazowych do standardowego rozkładu
normalnego (rozdział~\ref{ch:convergence}).
W szczególności, dla dowolnych rzeczywistych $a<b$,
\[
\lim_{n\to\infty}
\Prob\!\left(a<\frac{S_n-nm}{\sigma\sqrt n}\leq b\right)
=\Phi(b)-\Phi(a),
\]
gdzie $\Phi$ jest dystrybuantą standardowego rozkładu
normalnego.
\end{twierdzenie}
\begin{proof}
\emph{Krok pierwszy: sprowadzamy problem do średniej
zero i wariancji jeden.}
Zdefiniujmy $Z_k=(X_k-m)/\sigma$.
Funkcja $g(x)=(x-m)/\sigma$ jest ciągła, więc
$Z_k=g(X_k)$ są mierzalnymi zmiennymi losowymi.
Mają ten sam rozkład, równy obrazowi wspólnego
rozkładu $X_k$ przez $g$. Są wzajemnie niezależne,
gdyż dla różnych indeksów $k_1,\ldots,k_r$
i zbiorów borelowskich $B_1,\ldots,B_r$
\[
\begin{aligned}
\Prob\left(\bigcap_{j=1}^r\{Z_{k_j}\in B_j\}\right)
&=\Prob\left(\bigcap_{j=1}^r\{X_{k_j}\in g^{-1}(B_j)\}\right)\\
&=\prod_{j=1}^r\Prob(X_{k_j}\in g^{-1}(B_j))
=\prod_{j=1}^r\Prob(Z_{k_j}\in B_j).
\end{aligned}
\]
Przeciwobrazy $g^{-1}(B_j)$ są borelowskie,
więc można zastosować niezależność $X_{k_j}$.
Z algebry momentów (rozdział~\ref{ch:expectation}) mamy
\[
\E Z_k=\frac{\E X_k-m}{\sigma}=0,\qquad
\E Z_k^2=\frac{\E(X_k-m)^2}{\sigma^2}=1.
\]
Badana zmienna ma więc postać
$T_n=(Z_1+\cdots+Z_n)/\sqrt n$.
Jest mierzalna jako skończona kombinacja liniowa
zmiennych mierzalnych. Z powyższych tożsamości
i wzoru na wariancję sumy niezależnych zmiennych
(rozdział~\ref{ch:lln}) mamy
$\E T_n=0$ oraz $\Var(T_n)=n^{-1}\sum_{k=1}^n\Var(Z_k)=1$.

\emph{Krok drugi: rozwijamy funkcję charakterystyczną
jednego składnika.} Niech $\varphi(t)=\E e^{itZ_1}$.
Ponieważ $\E Z_1^2=1$, z twierdzenia o własnościach
funkcji charakterystycznej (rozdział~\ref{ch:expectation})
wynika $\varphi\in C^2(\R,\mathbb C)$ oraz
$\varphi(0)=1$, $\varphi'(0)=i\E Z_1=0$,
$\varphi''(0)=-\E Z_1^2=-1$.
Dla ustalonego $t\in\R$ połóżmy $h(s)=\varphi(st)$,
$s\in[0,1]$. Z reguły łańcuchowej mamy
$h'(s)=t\varphi'(st)$ oraz $h''(s)=t^2\varphi''(st)$.
Całkowanie przez części, zastosowane do części
rzeczywistej i urojonej, daje
\[
\int_0^1(1-s)h''(s)\,\dd\lambda_1(s)
=-h'(0)+h(1)-h(0).
\]
W ostatnim kroku korzystamy z równości
$\int_0^1h'(s)\,\dd\lambda_1(s)=h(1)-h(0)$,
która wynika z podstawowego twierdzenia rachunku
całkowego zastosowanego do funkcji $h$.
Po podstawieniu
otrzymujemy
\[
\varphi(t)=\varphi(0)+t\varphi'(0)
+t^2\int_0^1(1-s)\varphi''(st)\,\dd\lambda_1(s).
\]
Ponieważ $\int_0^1(1-s)\,\dd\lambda_1(s)=1/2$,
$\varphi(0)=1$, $\varphi'(0)=0$ oraz $\varphi''(0)=-1$,
otrzymujemy
\[
\begin{aligned}
\left|\varphi(t)-1+\frac{t^2}{2}\right|
&\leq t^2\int_0^1(1-s)|\varphi''(st)-\varphi''(0)|
\,\dd\lambda_1(s)\\
&\leq\frac{t^2}{2}
\sup_{|u|\leq |t|}|\varphi''(u)-\varphi''(0)|.
\end{aligned}
\]
Ciągłość $\varphi''$ w zerze oznacza, że ostatnie
supremum dąży do zera przy $t\to0$. Dostajemy zatem
\[
\varphi(t)=1-\frac{t^2}{2}+o(t^2)\qquad(t\to0).
\]
W dowodzie nie założyliśmy istnienia trzeciego
momentu: wystarcza skończony drugi moment.

\emph{Krok trzeci: wykorzystujemy niezależność.}
Z własności funkcji charakterystycznych
(rozdział~\ref{ch:expectation}) wynika
$\varphi_{aX}(t)=\varphi_X(at)$ oraz faktoryzacja
funkcji charakterystycznej sumy wzajemnie niezależnych
składników. Tę drugą własność stosujemy indukcyjnie:
suma pierwszych $k$ zmiennych jest niezależna od
$(k+1)$-szej, ponieważ wspólny rozkład wektora
jest iloczynem rozkładów i sumowanie jest borelowskie
(porównaj dowód nierówności maksymalnej Kołmogorowa
w rozdziale~\ref{ch:lln}).
Dla ustalonego $t\in\R$ mamy zatem
\[
\varphi_{T_n}(t)
=\left[\varphi\!\left(\frac{t}{\sqrt n}\right)\right]^n
=\left[1-\frac{t^2}{2n}+o\!\left(\frac1n\right)\right]^n.
\]
Dokładniej, dla ustalonego $t\ne0$ i
$r(h)=\varphi(h)-1+h^2/2$ mamy
$n r(t/\sqrt n)=t^2 r(t/\sqrt n)/(t/\sqrt n)^2\to0$.
Dla $t=0$ funkcje charakterystyczne mają wartość $1$.
Dla każdego ustalonego $t\in\R$ połóżmy
$u_n=n(\varphi(t/\sqrt n)-1)$.
W obu przypadkach mamy $u_n\to-t^2/2$.
Uzasadnijmy granicę potęg również dla zespolonych $u_n$.
Istnieje $M>0$ takie, że $|u_n|\leq M$ dla każdego $n$.
Dla $z\in\mathbb C$ szereg wykładniczy daje
\[
|e^z-1-z|\leq\frac{|z|^2}{2}e^{|z|}.
\]
Nierówność wynika z oszacowania
$1/j!\leq1/(2(j-2)!)$ dla $j\geq2$.
Dla $a_n=1+u_n/n$ i $b_n=e^{u_n/n}$ mamy
$|a_n|,|b_n|\leq e^{M/n}$ oraz
$|a_n-b_n|\leq M^2e^{M/n}/(2n^2)$.
Z algebraicznej tożsamości
$a^n-b^n=(a-b)\sum_{j=0}^{n-1}a^{n-1-j}b^j$
wynika
\[
|(1+u_n/n)^n-e^{u_n}|
\leq n e^{M(n-1)/n}\frac{M^2e^{M/n}}{2n^2}
=\frac{M^2e^M}{2n}\longrightarrow0.
\]
Ciągłość funkcji wykładniczej daje
$e^{u_n}\to e^{-t^2/2}$, a więc
$\varphi_{T_n}(t)\to e^{-t^2/2}$ dla każdego $t\in\R$.

\emph{Krok czwarty: identyfikujemy rozkład graniczny.}
W rozdziale~\ref{ch:expectation} wykazaliśmy, że
$e^{-t^2/2}$ jest funkcją charakterystyczną
$\Normal(0,1)$. Jest ona ciągła w zerze.
Z twierdzenia Lévy'ego o ciągłości
(rozdział~\ref{ch:convergence}) dla ciągu miar
$\Law(T_n)$ i jednoznaczności funkcji charakterystycznej
otrzymujemy $T_n\xrightarrow d\Normal(0,1)$.
Niech $F_n(x)=\Prob(T_n\leq x)$.
Dystrybuanta $\Phi$ rozkładu normalnego jest ciągła
na całej prostej, ponieważ rozkład ten ma gęstość
względem miary Lebesgue'a, a więc nie przypisuje
dodatniej masy żadnemu singletonowi
(rozdział~\ref{ch:random-variables}).
Z definicji zbieżności według rozkładu mamy więc
$F_n(x)\to\Phi(x)$ dla każdego $x\in\R$.
Z rozłącznego rozkładu zdarzenia $\{T_n\leq b\}$
na $\{T_n\leq a\}$ i $\{a<T_n\leq b\}$ otrzymujemy
\[
\Prob(a<T_n\leq b)=F_n(b)-F_n(a)
\longrightarrow\Phi(b)-\Phi(a).
\]
\end{proof}

\begin{uwaga}
CTG nie mówi, że $\overline X_n$ ma dla dużych $n$
dokładnie rozkład normalny. Mówi o granicy rozkładów
zmiennych \emph{wycentrowanych i wyskalowanych}.
Dla małych $n$ przybliżenie może być niedokładne,
zwłaszcza gdy rozkład $X_1$ jest silnie asymetryczny
lub ciężkoogonowy. Samo twierdzenie nie podaje
ilościowego oszacowania błędu przybliżenia.
Warunek $\sigma>0$ jest konieczny do standaryzacji:
jeśli $\Var(X_1)=0$, to z jednakowości rozkładów
oraz charakterystyki zmiennych o zerowej wariancji
(rozdział~\ref{ch:expectation}) wynika, że $X_k=m$
prawie na pewno dla każdego $k$ (wspólnie poza przeliczalną sumą
zbiorów zerowych), a więc $\overline X_n=m$
prawie na pewno dla każdego $n$.
\end{uwaga}

\section{Przykłady zastosowania CTG}

\begin{przyklad}[Przybliżenie de Moivre'a--Laplace'a]
Niech $B_1,B_2,\ldots$ będą wzajemnie niezależnymi
zmiennymi losowymi na jednej przestrzeni probabilistycznej,
każda o rozkładzie $\operatorname{Bern}(p)$ dla ustalonego
$p\in(0,1)$. Z twierdzenia o sumie prób Bernoulliego
(rozdział~\ref{ch:distributions}) wynika
$S_n=\sum_{k=1}^nB_k\sim\Bin(n,p)$.
Ponieważ $\E B_1=p$ i $\Var(B_1)=p(1-p)>0$,
wszystkie hipotezy klasycznego CTG są spełnione, więc
\[
\frac{S_n-np}{\sqrt{np(1-p)}}\xrightarrow d\Normal(0,1).
\]
Ponieważ mianownik $\sqrt{np(1-p)}$ jest dodatni,
dla dowolnych ustalonych $a<b$ zachodzi
\[
\begin{aligned}
&\Prob(np+a\sqrt{np(1-p)}<S_n\leq np+b\sqrt{np(1-p)})\\
&\qquad=\Prob\left(a<\frac{S_n-np}{\sqrt{np(1-p)}}\leq b\right)
\longrightarrow\Phi(b)-\Phi(a).
\end{aligned}
\]
Przybliżanie prawdopodobieństw zdarzeń o całkowitych
progach z poprawką ciągłości omówimy dalej.
Samo CTG nie podaje asymptotycznego równoważnika
prawdopodobieństwa pojedynczej wartości $S_n$.
Przypadki $p=0$ i $p=1$ są zdegenerowane: odpowiednio
$S_n=0$ lub $S_n=n$ prawie na pewno i standaryzacja
przez $\sqrt{np(1-p)}$ nie jest określona.
\end{przyklad}

\begin{przyklad}[Przybliżenie normalne rozkładu Poissona]
Dla każdej liczby rzeczywistej $\lambda>0$ niech
$N_\lambda\sim\Poi(\lambda)$; zmienne te nie muszą
być określone na wspólnej przestrzeni probabilistycznej.
Z rozdziału~\ref{ch:distributions} mamy
$\E N_\lambda=\Var(N_\lambda)=\lambda$.
Pokażemy, że przy $\lambda\to\infty$ zachodzi
$(N_\lambda-\lambda)/\sqrt\lambda\xrightarrow d\Normal(0,1)$.
Ze wzoru \eqref{eq:lotus} dla ograniczonej zespolonej
funkcji borelowskiej $x\mapsto e^{itx}$ i wzoru
na masy rozkładu Poissona otrzymujemy dla $t\in\R$
\[
\varphi_{N_\lambda}(t)=
e^{-\lambda}\sum_{k=0}^\infty\frac{(\lambda e^{it})^k}{k!}
=\exp(\lambda(e^{it}-1)).
\]
Szereg jest bezwzględnie zbieżny, gdyż suma modułów
po uwzględnieniu czynnika $e^{-\lambda}$ wynosi jeden.
Z własności funkcji charakterystycznej przekształcenia
afinicznego (rozdział~\ref{ch:expectation}) otrzymujemy
\[
\varphi_{(N_\lambda-\lambda)/\sqrt\lambda}(t)
=\exp\!\left(\lambda\left[e^{it/\sqrt\lambda}
-1-\frac{it}{\sqrt\lambda}\right]\right).
\]
Z bezwzględnie zbieżnego szeregu wykładniczego
i nierówności $1/j!\leq 1/(6(j-3)!)$ dla $j\geq3$
wynika dla $z\in\mathbb C$
\[
|e^z-1-z-z^2/2|\leq |z|^3 e^{|z|}/6.
\]
Podstawiając $z=it/\sqrt\lambda$, dostajemy
\[
\left|\lambda\left(e^{it/\sqrt\lambda}-1-
\frac{it}{\sqrt\lambda}\right)+\frac{t^2}{2}\right|
\leq\frac{|t|^3}{6\sqrt\lambda}e^{|t|/\sqrt\lambda}
\longrightarrow0
\]
dla każdego ustalonego $t\in\R$.
Z ciągłości funkcji wykładniczej funkcje charakterystyczne
standaryzowanych zmiennych dążą punktowo do
$e^{-t^2/2}$, funkcji charakterystycznej $\Normal(0,1)$.
Dla dowolnego ciągu $\lambda_j\to\infty$ twierdzenie
Lévy'ego (rozdział~\ref{ch:convergence}) daje zbieżność
dystrybuant standaryzowanych zmiennych w każdym
$x\in\R$, ponieważ $\Phi$ jest ciągła.
Gdyby dla pewnego $x$ granica dystrybuant przy
rzeczywistym parametrze $\lambda\to\infty$ nie była
równa $\Phi(x)$, istniałby ciąg $\lambda_j\geq j$
i $\varepsilon>0$, dla których moduł różnicy byłby
co najmniej $\varepsilon$. Przeczy to uzyskanej
zbieżności dla każdego takiego ciągu.
\end{przyklad}

\begin{przyklad}[Od rozkładu Studenta do normalnego]
Niech $Z,W_1,W_2,\ldots$ będą wzajemnie niezależnymi
zmiennymi o rozkładzie $\Normal(0,1)$ na wspólnej
przestrzeni probabilistycznej.
Dla całkowitego $n\geq1$ połóżmy $U_n=\sum_{k=1}^nW_k^2$.
Z własności rozkładu chi-kwadrat
(rozdział~\ref{ch:distributions}) mamy $U_n\sim\chi_n^2$.
Zmienna $U_n$ jest nieujemna i mierzalna jako
skończona suma kwadratów zmiennych mierzalnych.
Ponieważ $\{U_n=0\}\subseteq\{W_1=0\}$, a rozkład
$W_1$ ma gęstość, zachodzi $\Prob(U_n=0)=0$.
Zmienna $Z$ jest niezależna od $U_n$: rozkład łączny
$Z$ i wektora $(W_1,\ldots,W_n)$ jest iloczynem
ich rozkładów, a suma kwadratów jest funkcją ciągłą
tego wektora.
Definiujemy
\[
T_n=\begin{cases}
Z/\sqrt{U_n/n},&U_n>0,\\
0,&U_n=0.
\end{cases}
\]
Jest to mierzalna zmienna losowa i z definicji
rozkładu Studenta (rozdział~\ref{ch:distributions})
mamy $T_n\sim t_n$.
Zmienne $W_k^2$ są wzajemnie niezależne i jednakowo
rozłożone, ponieważ powstają przez zastosowanie
tej samej ciągłej funkcji $x\mapsto x^2$ do
niezależnych zmiennych o wspólnym rozkładzie.
Są całkowalne, gdyż $\E W_k^2=1$.
Z MPWL Kołmogorowa (rozdział~\ref{ch:lln}) istnieje
mierzalny zbiór $G$ o prawdopodobieństwie jeden,
na którym $U_n/n\to1$.
Po usunięciu z $G$ przeliczalnej sumy zdarzeń
$\{U_n=0\}$ nadal mamy zbiór pełnej miary.
Na tym zbiorze ciągłość funkcji
$x\mapsto1/\sqrt x$ na $(0,\infty)$ daje
\[
T_n(\omega)=Z(\omega)(U_n(\omega)/n)^{-1/2}
\longrightarrow Z(\omega).
\]
Otrzymujemy zatem zbieżność prawie na pewno w tej
konstrukcji, a więc także według rozkładu
(rozdział~\ref{ch:convergence}).
Alternatywnie $\sqrt{U_n/n}\to1$ według
prawdopodobieństwa i $Z\xrightarrow{d}Z$ dla stałego
ciągu zmiennych, więc twierdzenie Slutskiego o ilorazie
daje ten sam wniosek o rozkładach.
W ten sposób uzasadniamy granicę rozkładu Studenta
przy całkowitej liczbie stopni swobody dążącej
do nieskończoności. W rozdziale~\ref{ch:distributions}
rozważono także granicę dla rzeczywistego parametru.
\end{przyklad}

\section{Co rozróżniają prawa graniczne?}

Niech $X_1,X_2,\ldots$ będą wzajemnie niezależnymi
rzeczywistymi zmiennymi losowymi na jednej przestrzeni
probabilistycznej, o tym samym rozkładzie i
$\E|X_1|<\infty$. Oznaczmy $m=\E X_1$.
Z rozdziału~\ref{ch:lln} wiemy, że
$\overline X_n\xrightarrow{\Prob}m$ (SPWL) oraz
$\overline X_n\to m$ prawie na pewno (MPWL).
Oba prawa dotyczą zbieżności średnich do stałej.
Klasyczne CTG wymaga ponadto $0<\Var(X_1)<\infty$.
Dla $\sigma=\sqrt{\Var(X_1)}$ opisuje ono zbieżność
rozkładów zmiennych
$T_n=\sqrt n(\overline X_n-m)/\sigma$
do niezdegenerowanego rozkładu normalnego.
Dla każdego ustalonego $c>0$ mamy
\[
\begin{aligned}
\Prob\left(m-\frac{c\sigma}{\sqrt n}<\overline X_n
\leq m+\frac{c\sigma}{\sqrt n}\right)
&=\Prob(-c<T_n\leq c)\\
&\longrightarrow\Phi(c)-\Phi(-c)=2\Phi(c)-1.
\end{aligned}
\]
Ostatnia równość wynika z symetrii gęstości
standardowego rozkładu normalnego.
W szczególności dla $c=2$ otrzymujemy graniczne
prawdopodobieństwo $\Phi(2)-\Phi(-2)\approx0{,}9545$.
W przeciwieństwie do prawa wielkich liczb jest to
informacja o prawdopodobieństwie przynależności średniej
do przedziału, którego długość maleje jak $n^{-1/2}$.

Na tym kończy się pierwszy blok twierdzeń granicznych.
W kolejnych rozdziałach uzupełnimy go o warunkową wartość
oczekiwaną, koncentrację miary i granice dla
niejednakowo rozłożonych składników; powrócimy też
do problemów przewidywania i empirycznego opisu rozkładu.

\begin{twierdzenie}[Jednostajna zbieżność dystrybuant do granicy ciągłej]
Niech $(F_n)_{n\geq1}$ będzie ciągiem dystrybuant
na $\R$, a $F$ niech będzie \emph{ciągłą} dystrybuantą.
Jeżeli $F_n(x)\to F(x)$ dla każdego $x\in\R$,
to
\[
\lim_{n\to\infty}\sup_{x\in\R}|F_n(x)-F(x)|=0.
\]
\end{twierdzenie}
\begin{proof}
Ustalmy $\varepsilon>0$. Z własności dystrybuant
(rozdział~\ref{ch:random-variables}) mamy
$F(x)\to0$ przy $x\to-\infty$ i $F(x)\to1$
przy $x\to+\infty$. Wybierzmy więc $a<b$ tak,
by $F(a)<\varepsilon$ oraz $1-F(b)<\varepsilon$.
Z jednostajnej ciągłości $F$ na zwartym przedziale
$[a,b]$ istnieje $\delta>0$ takie, że
$|F(x)-F(y)|<\varepsilon$ dla $x,y\in[a,b]$
spełniających $|x-y|<\delta$.
Wybierzmy całkowite $m>(b-a)/\delta$ i połóżmy
$x_j=a+j(b-a)/m$ dla $j=0,\ldots,m$.
Wówczas $0\leq F(x_{j+1})-F(x_j)<\varepsilon$
dla $j=0,\ldots,m-1$.
Dla każdego z $m+1$ punktów podziału, z założonej
zbieżności punktowej, istnieje $N_j$ takie, że
$|F_n(x_j)-F(x_j)|<\varepsilon$ dla $n\geq N_j$.
Przyjmując $N=\max_{0\leq j\leq m}N_j$,
otrzymujemy wszystkie te nierówności jednocześnie
dla $n\geq N$.
Niech $n\geq N$. Jeśli $x_j\leq x\leq x_{j+1}$,
to z monotoniczności obu dystrybuant dostajemy
\[
\begin{aligned}
F_n(x)-F(x)
&\leq F_n(x_{j+1})-F(x_j)\\
&=F_n(x_{j+1})-F(x_{j+1})
  +F(x_{j+1})-F(x_j)<2\varepsilon,\\
F(x)-F_n(x)
&\leq F(x_{j+1})-F_n(x_j)\\
&=F(x_{j+1})-F(x_j)+F(x_j)-F_n(x_j)
<2\varepsilon.
\end{aligned}
\]
Jeżeli $x<a$, to
$0\leq F(x)\leq F(a)<\varepsilon$ oraz
$0\leq F_n(x)\leq F_n(a)<F(a)+\varepsilon<2\varepsilon$.
Stąd $|F_n(x)-F(x)|<2\varepsilon$.
Jeżeli $x>b$, to
$0\leq1-F(x)\leq1-F(b)<\varepsilon$ oraz
$0\leq1-F_n(x)\leq1-F_n(b)<2\varepsilon$,
więc także $|F_n(x)-F(x)|<2\varepsilon$.
Łącząc wszystkie przypadki, otrzymujemy dla $n\geq N$
\[
\sup_{x\in\R}|F_n(x)-F(x)|\leq2\varepsilon.
\]
Indeks $N$ zależy od $\varepsilon$, lecz nie od $x$.
Dowolność $\varepsilon>0$ dowodzi zbieżności jednostajnej.
\end{proof}

\begin{uwaga}
W szczególności dla standaryzowanych sum $T_n$
z klasycznego CTG, ponieważ dystrybuanta normalna
$\Phi$ jest ciągła, zbieżność według rozkładu daje
zbieżność dystrybuant w każdym punkcie. Na mocy
powyższego twierdzenia
\[
d_n:=\sup_{x\in\R}|\Prob(T_n\leq x)-\Phi(x)|
\longrightarrow0.
\]
Dla dowolnych rzeczywistych $a<b$ mamy zatem
\[
|\Prob(a<T_n\leq b)-(\Phi(b)-\Phi(a))|
\leq 2d_n.
\]
Istotnie, zapisujemy prawdopodobieństwo jako różnicę
dystrybuant i stosujemy nierówność trójkąta.
Oszacowanie jest jednostajne względem obu końców
przedziału, więc obowiązuje również dla progów
zależnych od $n$. Jednostajna zbieżność dystrybuant
nie jest tu dodatkowym założeniem, lecz konsekwencją
zbieżności według rozkładu do ciągłej dystrybuanty.
Nadal nie otrzymujemy ilościowego oszacowania
szybkości zbieżności $d_n\to0$.
\end{uwaga}

\section{Precyzja przybliżenia normalnego i dalsze przykłady}

\begin{przyklad}[Odchylenia losowego spaceru]
Niech $\xi_1,\xi_2,\ldots$ będą wzajemnie niezależnymi
rzeczywistymi zmiennymi losowymi na jednej przestrzeni
probabilistycznej, o jednakowym rozkładzie zadanym przez
$\Prob(\xi_k=1)=\Prob(\xi_k=-1)=1/2$.
Niech $S_n=\sum_{k=1}^n\xi_k$.
Z definicji wartości oczekiwanej i wariancji mamy
$\E\xi_k=(1-1)/2=0$ oraz
$\Var(\xi_k)=\E\xi_k^2=1$.
Klasyczne CTG daje $S_n/\sqrt n\xrightarrow d\Normal(0,1)$.
Dla $a>0$ zbiór $[-a,a]$ ma brzeg $\{-a,a\}$,
który ma miarę zero względem standardowego rozkładu
normalnego. Z \atlasref{ME-22::thm:portmanteau-22}{twierdzenia Portmanteau}
w wersji dla zbiorów ciągłości miary granicznej
(porównaj rozdział~\ref{ch:convergence}) oraz
symetrii dystrybuanty normalnej wynika więc
\[
\Prob(|S_n|\leq a\sqrt n)
\longrightarrow\Phi(a)-\Phi(-a)=2\Phi(a)-1.
\]
Uwzględnienie obu końców domkniętego przedziału
nie zmienia granicy, choć może zmieniać
prawdopodobieństwo dla ustalonego $n$.
Ponieważ $\E|\xi_1|=1$, MPWL Kołmogorowa
(rozdział~\ref{ch:lln}) daje także $S_n/n\to0$
prawie na pewno. Skale $n$ i $\sqrt n$ odpowiadają
dwóm różnym pytaniom: zanikowi średniego
odchylenia i rozkładowi wahań po przeskalowaniu.
\end{przyklad}

\begin{przyklad}[Przybliżenie dwumianowe z poprawką ciągłości]
Niech $n\geq1$ będzie całkowite, niech
$S_n\sim\Bin(n,p)$ dla ustalonego $p\in(0,1)$
oraz oznaczmy $\mu=np$, $s=\sqrt{np(1-p)}>0$.
Dla liczb całkowitych $a\leq b$, ponieważ $S_n$
przyjmuje wartości całkowite prawie na pewno, mamy
dokładną równość prawdopodobieństw
\[
\Prob(a\leq S_n\leq b)
=\Prob\left(\frac{a-1/2-\mu}{s}
<\frac{S_n-\mu}{s}\leq\frac{b+1/2-\mu}{s}\right).
\]
Przybliżając dystrybuantę standaryzowanego $S_n$
dystrybuantą $\Phi$, otrzymujemy
\[
\Prob(a\leq S_n\leq b)\approx
\Phi\!\left(\frac{b+1/2-\mu}{s}\right)-
\Phi\!\left(\frac{a-1/2-\mu}{s}\right).
\]
Przesunięcie progów o $1/2$ odpowiada zastąpieniu
każdej wartości całkowitej $k$ przedziałem
$(k-1/2,k+1/2]$ o długości jeden.
Niech
\[
d_n=\sup_{x\in\R}
\left|\Prob\left(\frac{S_n-\mu}{s}\leq x\right)-\Phi(x)\right|.
\]
Z poprzedniej sekcji wiadomo, że $d_n\to0$
dla ustalonego $p\in(0,1)$. Na podstawie dokładnej
równości zdarzeń z dokładnością do zbioru zerowego
i oszacowania różnicy dystrybuant błąd bezwzględny
powyższego przybliżenia nie przekracza $2d_n$.
Oszacowanie obowiązuje również dla całkowitych
progów $a,b$ zależnych od $n$. Nie oznacza ono,
że poprawka ciągłości zawsze zmniejsza błąd
w porównaniu z innym przybliżeniem, ani że błąd
względny małych prawdopodobieństw dąży do zera.
Samo CTG nie określa szybkości zbieżności $d_n$.
\end{przyklad}

\begin{przyklad}[Gdy klasyczne CTG nie ma zastosowania]
Niech $C_1,C_2,\ldots$ będą wzajemnie niezależnymi
zmiennymi losowymi na jednej przestrzeni
probabilistycznej, o standardowym rozkładzie Cauchy'ego.
Jak wykazano w rozdziale~\ref{ch:lln}, korzystając
z funkcji charakterystycznych, dla każdego $n\geq1$
zachodzi
\[
\Law\left(\frac{C_1+\cdots+C_n}{n}\right)=\Law(C_1).
\]
Rozkład Cauchy'ego nie ma skończonego pierwszego
momentu bezwzględnego ani skończonego drugiego
momentu (rozdział~\ref{ch:distributions}).
Średnie mają stale ten sam niezdegenerowany
rozkład, więc z jednoznaczności granicy według
rozkładu (rozdział~\ref{ch:convergence}) nie mogą
zbiegać według rozkładu do stałej.
W szczególności nie są określone parametry $m$
i $\sigma$ potrzebne do standaryzacji w klasycznym
CTG. Przykład nie wyklucza innych twierdzeń
granicznych przy innych normalizacjach.
\end{przyklad}

\begin{uwaga}[Gdzie w dowodzie CTG używa się drugiego momentu?]
Niech $Z$ będzie rzeczywistą zmienną losową,
$\E Z=0$ i $\E Z^2=1$.
Z nierówności Cauchy'ego--Schwarza wynika
$\E|Z|\leq(\E Z^2)^{1/2}=1$.
Z twierdzenia o własnościach funkcji charakterystycznej
(rozdział~\ref{ch:expectation}) otrzymujemy
\[
\varphi_Z'(t)=\E(iZe^{itZ}),\qquad
\varphi_Z''(t)=\E(-Z^2e^{itZ}).
\]
W uzasadnieniu różniczkowania używa się dla
rzeczywistego $h\ne0$ nierówności
$|e^{ihz}-1|/|h|\leq|z|$ dla $z\in\R$,
wynikającej z $|e^{iu}-1|=2|\sin(u/2)|\leq|u|$.
Ilorazy różnicowe pod całką są więc zmajoryzowane
odpowiednio przez całkowalne funkcje $|Z|$ i $Z^2$.
Ciągłość drugiej pochodnej wynika również z
\atlasref{ME-08::thm:dct}{twierdzenia o zbieżności zmajoryzowanej}:
dla $t_j\to t$ funkcje $-Z^2e^{it_jZ}$ zbiegają
punktowo do $-Z^2e^{itZ}$, a ich moduły są równe $Z^2$.
Twierdzenie stosujemy do części rzeczywistej
i urojonej. Rozwinięcie Taylora z resztą Peana,
uzasadnione w dowodzie klasycznego CTG, daje
$\varphi_Z(t)=1-t^2/2+o(t^2)$.
Nie potrzeba założenia $\E|Z|^3<\infty$.
Ilościowe oszacowanie błędu przybliżenia normalnego
wymaga odrębnego twierdzenia; nie wynika z samej
zbieżności funkcji charakterystycznych wykazanej
powyżej. W rozdziale~\ref{ch:advancedlimits}
omawiamy natomiast między innymi nierówności
koncentracji dla średnich oraz CTG dla składników
o niejednakowych rozkładach. Nierówności koncentracji
oszacowują prawdopodobieństwa odchyleń, a nie
bezpośrednio odległość dystrybuanty od $\Phi$.
\end{uwaga}

